% Volume~III, appendix B: tools (English version). Translator EN2, 2026-10-29, from
% memoire/tome3/fr/annexes/B-outils.tex (writer D); same labels, same macros, same numbers.
% Model: memoire/tome2/en/annexes/B-outils.tex. The tables of the tools and of the figures are generated by
% memoire/outils/tome3_annexes.py from commit 4b071d0 (memoire/tome3/annexes/outils-tome3.tex, figures-tome3.tex,
% bilingual by \EF).

\chapter{Tools}\label{t3ann:outils}

\begin{origine}
\textbf{At the start}, the appendix on the tools of Volumes~I and II: each tool sorted by study or by lot, with the first sentence of its own description, and the chain that builds the thesis.
\textbf{This appendix establishes} the same for Volume~III: the chain that builds it and what it checks; each committed tool of the quest, of mimicry, of the leads, of the judges of the figures and of the thesis, sorted by chapter; each figure in the style of the volume and the files it reads; the licences of the dependencies.
\textbf{What remains} \barre{is to put the tool of the appendices into the tome3 target of construire.sh}: nothing, for the chain; the tool of the appendices has been there since 2026-10-29 (commit 57db4ce), and it reads the commit rather than the working tree since commit e5ce266 (\cref{t3ann:sec:chaine}).
\end{origine}

\section{The chain of Volume~III}\label{t3ann:sec:chaine}

Volume~III is built like the first two, by the same script and the same \texttt{latexmk} configuration (\fichier{memoire/construire.sh}, \fichier{memoire/latexmkrc}), from \fichier{memoire/}:
\begin{itemize}
\item \texttt{./construire.sh tome3} (or \texttt{tome3-fr}, \texttt{tome3-en}) first runs the eight data tools of the figures, one per lot, under \fichier{memoire/outils/}: \texttt{tome3\_} followed by \texttt{eclaireur}, \texttt{genome}, \texttt{duplication}, \texttt{synthese}, \texttt{codon}, \texttt{synonyme}, \texttt{juge} or \texttt{latents}, then \texttt{.py}. Each rereads the committed logs, writes its CSVs and its numbers under \fichier{memoire/figures/tome3/donnees/}, and stops at its first guard; their regeneration gives the committed files again to the bit (\fichier{memoire/tome3/JUGEMENT-figures3.md}, § 4.3).
\item Next, it generates the four tables of appendices A, B and C by \fichier{memoire/outils/tome3_annexes.py} (\cref{t3ann:tab:seuils-dates,t3ann:tab:outils,t3ann:tab:figures,t3ann:tab:journaux}), from the \barre{commit; it refuses to run if a file it counts differs from the commit} publication commit of the volume, \texttt{4b071d0} (paragraph ``The tool of the appendices'', below).
\item It then compiles, separately, the figures of the spectra lead: \fichier{genome/pistes/spectres/figures/spectres-figures.tex} is a standalone document whose data paths are relative to its folder, which is read-only; it is therefore compiled in that folder, with its output under \fichier{memoire/tome3/pistes/} (ignored by git), where the chapter on the spectra takes its pages.
\item It finally builds the two PDFs by \texttt{latexmk} (lualatex, biber) in \fichier{memoire/tome3/fr/} and \fichier{memoire/tome3/en/}, and stops with an error on any undefined reference or citation and on any missing data file.
\end{itemize}
Two tools have no place in a build, and are not in it: \fichier{memoire/outils/tome3_palette.py}, which rewrites its files with other line endings (equal in content only), and \fichier{memoire/outils/tome3_cassures.py}, which forces the guards to false (appendix~\ref{t3ann:protocoles}).

\paragraph{The tool of the appendices.} \fichier{memoire/outils/tome3_annexes.py} generates the four tables of appendices A, B and C from the \barre{HEAD commit} publication commit of the volume (\texttt{git ls-tree}, \texttt{git cat-file}, \texttt{git log}): the tools, the figures and their data, the log folders, the dated threshold files. Its control: every file cited by the ``Data'' row of a figure (full path, patterns \texttt{*} and \texttt{\{a,b\}} allowed) exists in the repository; its break (\texttt{--casser}, a non-existent file added to the first figure) makes it fall, code 1 \mes. \barre{It is not yet in the tome3 target of construire.sh, which this writer was not allowed to modify}; it has been there since 2026-10-29 (commit 57db4ce): the \texttt{tome3} target of \fichier{memoire/construire.sh} runs it before latexmk, and each build of the volume rewrites its tables from the repository. Corrected on 2026-10-29 (commit e5ce266, grave G3 of \fichier{memoire/tome3/JUGEMENT-final.md}): it counted the lines and the bytes in the working tree, so that an uncommitted change from another session gave there 124 lines to \fichier{outils/mimetisme/humain/e4b_defauts.py}, which has 122 in the commit, and the sizes of appendix~\ref{t3ann:journaux} depended on the line endings of the copy (\fichier{mesures/genome/construction/codon/}: \SI{2.47}{\Mo} in one copy, \SI{2.43}{\Mo} in the commit). It now counts the lines on the content of the git objects and the sizes on their size, \barre{and it refuses to run, code 1, if a file it counts differs from the commit in the working tree} \mes. \emph{Pinned on 2026-10-29} to the publication commit of the volume, \texttt{4b071d0}, like the tools of Volumes~I and II: read at the HEAD commit, it refused to run as soon as another session left a counted file uncommitted; read at this commit, it gives back the four tables bit for bit, and neither the working tree nor later commits enter them any more.

\section{The committed tools}\label{t3ann:sec:outils}

Every number \mes{} or \calc{} of this volume comes out of a committed tool. \Cref{t3ann:tab:outils} lists them, sorted by chapter, each with the first line of its own description, quoted as written (in French), and its number of lines. The tools of the quest, of mimicry and of the leads belong to their streams; the thesis cites them and does not modify them. Those of the second draw of P3, in progress in another session, are not in it (reserve R6).

\input{\Racine/tome3/annexes/outils-tome3}

\section{Each figure and the files it reads}\label{t3ann:sec:figures}

A beautiful figure remains a true figure: it is drawn from the committed data, by a committed script. \Cref{t3ann:tab:figures} gives, for each figure in the style of Volume~III, the files that its ``Data'' row cites; they are the ones read by the tool of the lot named at the end of the row. The figures outside the common style (mimicry, the two leads) cite their data in their caption, in the chapter that places them.

\input{\Racine/tome3/annexes/figures-tome3}

\section{The licences}\label{t3ann:sec:licences}

The SOCLE asks for the licence of each dependency and of each studied repository to be written down. No external code repository was read by the quest for the genome: the state of the art read only articles and their abstracts (\fichier{genome/00-etat-de-l-art.md}), and no code licence is at stake there. The dependencies are those of \cref{t3ann:tab:licences}, each read in the package installed on the machine.

\begin{table}[ht]
\centering\footnotesize
\caption[The licences of the dependencies of Volume~III]{The licences of the dependencies of the tools and of the thesis of Volume~III, read in the packages installed on 29 October 2026.}
\label{t3ann:tab:licences}
\begin{tabularx}{\linewidth}{@{}>{\raggedright\arraybackslash}p{0.22\linewidth}>{\raggedright\arraybackslash}p{0.30\linewidth}>{\raggedright\arraybackslash}X@{}}
\toprule
Dependency & Licence & Where it is read \\
\midrule
numpy 2.5.4 & BSD-3-Clause (and 0BSD, MIT, Zlib, CC0-1.0 for parts) & \texttt{License-Expression} of \texttt{numpy-2.5.4.dist-info/METADATA} \\
scipy 1.18.1 & three-clause BSD & \texttt{scipy-1.18.1.dist-info/LICENSE.txt} \\
rayon 1.12.0, half 2.7.1, anyhow 1.0.104 (Rust) & MIT or Apache-2.0 & \texttt{license} field of their \texttt{Cargo.toml} (local Cargo registry) \\
llama.cpp b11515 & MIT & Volume~I, appendix on the tools (official sources of the same commit) \\
pgfplots & GNU GPL, version 3 or later & header of \texttt{pgfplots.code.tex} (MiKTeX) \\
PGF/TikZ & LPPL and/or GNU GPL & header of \texttt{tikz.sty} (MiKTeX) \\
biblatex & LPPL 1.3c & header of \texttt{biblatex.sty} (MiKTeX) \\
\bottomrule
\end{tabularx}
\end{table}

\section*{What could be wrong}

\begin{faux}
\begin{itemize}
\item \textbf{The first line of a description} is the one the tool gives of itself; it may have aged. The tool of the appendices quotes it, it does not judge it.
\item \textbf{The control of the figures} checks that the cited files exist, not that they are those the tool really reads; it is the bit-exact regeneration of the data (judge of the third pass of the figures) that checks this.
\item \textbf{The licences} are read in the installed packages; another version may have another one.
\end{itemize}
\end{faux}
